\chapter{Share Classes, Holding Companies, Closed-End Funds and SPACs}\label{ch:s1:share-classes-holding-companies-closed-end-funds-spacs}

A closed-end fund holds \$100 million of listed stocks and trades at \$85 million. Its shares cannot be redeemed for the stocks, so nothing forces the
two values together: the discount can last for years, narrow when sentiment turns, or widen. Buying the fund and selling its holdings locks in
fifteen cents on the dollar only if something closes the gap. On the synthetic model of this chapter the trade earns 0.87\% a year if the discount is left
to drift back toward its average, with a median holding of three years and a two-in-five chance of losing; with catalysts (a tender, a liquidation, an
activist) arriving at 0.3 a year it earns 5.0\% a year and closes in a median 1.4 years. Structural relative value is the family of trades in which a
contract or a structure ties two prices together: share classes, holding companies, closed-end funds and special-purpose acquisition companies. The
build is \texttt{firm.structrv}.

\section{Dual share classes}

\begin{definition}[Dual-class share]\label{def:s1:share-classes-holding-companies-closed-end-funds-spacs:dual}
A \emph{dual-class share}\index{dual-class share} structure gives a company two or more listed classes of common stock with the same claim on its cash flows but different rights
(usually votes), so that the classes should trade close to each other and their spread reflects the value of the difference in rights, liquidity and index
eligibility.
\end{definition}

Share classes are the cleanest structural pair: the same company, the same dividends, a spread that should be small and stationary. They are the real
counterpart of chapter~4's planted twins, whose log price spread reverted with a ten-day time constant and a 2\% standard deviation, and which the
\omterm{def:s1:pairs-and-baskets:distance}{distance method} found and traded at a Sharpe ratio of 5 after costs. The real spreads are driven by what differs between the classes: voting rights
matter in takeover contests, one class may be in an index and the other not (chapter~10), and one may be far less liquid. A share-class book trades the
spread against its own history and closes positions before events that could change the value of the vote.

\section{Holding-company discounts}

\begin{definition}[Holding-company discount]\label{def:s1:share-classes-holding-companies-closed-end-funds-spacs:holding}
A \emph{holding-company discount}\index{holding-company discount} is the amount by which a holding company's market value falls short of the value of its stakes in other listed and
unlisted companies (net of its debt), expressed as a fraction of that value.
\end{definition}

A holding company that owns listed stakes can be hedged exactly: buy the holding company, sell its stakes in proportion. What is left is the discount,
which has reasons to exist (taxes on a sale of the stakes, the holding's costs, governance) and few forces to close it; unlike a closed-end fund, a holding
company rarely liquidates. The chapter's holding company trades at 70 cents on the dollar with a discount reverting toward 80 cents over a half-life of three
years and no catalyst: the hedged trade earns 0.94\% a year, closes within three years in only 2\% of the paths, and loses in 40\%.

\section{Closed-end fund discounts}

\begin{definition}[Closed-end fund discount]\label{def:s1:share-classes-holding-companies-closed-end-funds-spacs:cef}
A \emph{closed-end fund discount}\index{closed-end fund discount} is the amount by which a closed-end fund's share price falls short of its net asset value per share, as a fraction of the
net asset value; a closed-end fund issues a fixed number of shares that trade on an exchange and cannot be redeemed at net asset value.
\end{definition}

The closed-end fund puzzle has two parts. Funds trade at discounts most of the time, and those discounts move. Lee, Shleifer and Thaler traced the movement to
individual investors' sentiment: discounts on different funds move together, new funds get started when seasoned ones sell near or above their asset value,
and discounts narrow when small stocks do well. Pontiff found that deviations are larger where arbitrage is costlier: portfolios hard to replicate, small
dividends, small funds and high interest rates, which together explained a quarter of the cross-section of mispricing.

\begin{center}
\small
\begin{tabular}{@{}lcccc@{}}
\toprule
& \multicolumn{3}{c}{closed-end fund at 85 cents} & holding company\\
catalysts a year & none & 0.1 & 0.3 & at 70 cents, none\\
\midrule
mean return per trade; a year & 1.9\%; 0.87\% & 4.4\%; 2.23\% & 7.9\%; 5.04\% & 2.8\%; 0.94\%\\
median years held & 2.96 & 2.13 & 1.37 & 3.00\\
closed within three years & 50\% & 63\% & 79\% & 2\%\\
closed by a catalyst & 0\% & 19\% & 46\% & 0\%\\
5th percentile of return; share that lose & $-14.7\%$; 40\% & $-13.3\%$; 30\% & $-10.4\%$; 17\% & $-13.9\%$; 40\%\\
\bottomrule
\end{tabular}
\end{center}

The model (\cref{lst:s1:share-classes-holding-companies-closed-end-funds-spacs:discount}) is a mean-reverting discount, in logs: the fund starts at 85 cents
on the dollar ($-16.3\%$) and reverts toward 90 cents ($-10.5\%$) with a half-life of two years and a volatility of 8\% a year. The trade buys the fund and
shorts its portfolio, pays 1.3\% a year in carry (the fund's expenses and the hedge's borrow), and closes when the discount reaches 5\% or after three years.
Without a catalyst, most of the expected narrowing is to the average discount, not to zero, and the carry eats much of it: the trade is a slow, uncertain
2\% over three years. A catalyst changes the arithmetic: the discount closes to 2\% at once, the capital is freed early, and the return per year rises
nearly sixfold at 0.3 catalysts a year (\cref{fig:s1:share-classes-holding-companies-closed-end-funds-spacs:exit}). This is why discount trades are so often
activist trades: the investor who can force a tender or a liquidation owns the catalyst.

\begin{omfigure}
\begin{tikzpicture}
\begin{axis}[ybar,bar width=7pt,width=11cm,height=5.0cm,xlabel={\small years until the trade closes},ylabel={\small share of paths (\%)},
  tick label style={font=\footnotesize},xmin=0,xmax=3.1,ymin=0,ymax=60,grid=major,grid style={gray!25},
  legend style={font=\scriptsize,at={(0.03,0.97)},anchor=north west},table/col sep=comma]
\addplot[fill=omThm!60,draw=omThm] table[x=years,y=no_catalyst]{figdata/strategies-1/12-share-classes-holding-companies-closed-end-funds-spacs/exit.csv};
\addplot[fill=omDef!60,draw=omDef] table[x=years,y=catalyst]{figdata/strategies-1/12-share-classes-holding-companies-closed-end-funds-spacs/exit.csv};
\legend{no catalyst,0.3 catalysts a year}
\end{axis}
\end{tikzpicture}

\omcaption{The \omterm{def:s1:share-classes-holding-companies-closed-end-funds-spacs:cef}{closed-end fund discount} trade: when it closes (the discount reaches 5\%, a catalyst arrives, or three years pass), over 10,000 simulated
paths; the last bar holds the trades still open at three years. Data: \texttt{s1\_structrv.fund}.}
\label{fig:s1:share-classes-holding-companies-closed-end-funds-spacs:exit}
\end{omfigure}

\section{Blank-cheque companies and their trust value}

\begin{definition}[Special purpose acquisition company, trust value]\label{def:s1:share-classes-holding-companies-closed-end-funds-spacs:spac}
A \emph{special purpose acquisition company}\index{special purpose acquisition company} (SPAC) is a listed shell company that raises cash in an initial public offering, holds it in a trust, and has a
limited time to merge with a private company; its shareholders can redeem their shares for their share of the trust instead of taking part in the merger.
The \emph{trust value}\index{trust value} per share is that redemption amount, the cash raised plus the interest the trust has earned.
\end{definition}

The redemption right makes a SPAC share a floor plus an option: at the vote the holder can take the \omterm{def:s1:share-classes-holding-companies-closed-end-funds-spacs:spac}{trust value} or keep a share of the merged company. A
share bought below the \omterm{def:s1:share-classes-holding-companies-closed-end-funds-spacs:spac}{trust value} has a guaranteed minimum return (if the trust is safe) and a free call on the merger. Bought at \$9.80 with a trust of
\$10.00 earning 4\% a year and the vote a year away, the floor is \$10.41, a 6.2\% return (\cref{lst:s1:share-classes-holding-companies-closed-end-funds-spacs:spac}).
What the call is worth depends on the merged company. Klausner, Ohlrogge and Ruan found that SPACs raising \$10.00 a share held far less net cash per share
at their mergers (a median of \$5.70 for those merging from January 2019 to June 2020), and that shareholders who held through the merger suffered steep
losses while sponsors profited. The chapter's assumption follows that record: a merged value lognormal around \$7 with a volatility of 60\%. The arbitrageur
who redeems when the merged value is below the trust (74\% of the paths) earns 20.2\% on average, the floor plus a call on a volatile outcome; a holder who
never redeems loses 13.9\% on average and 27.8\% at the median.

\begin{dated}{2026-09}{SEC rules for SPACs}\label{dat:s1:share-classes-holding-companies-closed-end-funds-spacs:sec}
On January 24, 2024 the SEC adopted rules and amendments for SPAC initial public offerings and de-SPAC transactions (the mergers), requiring among other
things enhanced disclosure about conflicts of interest, sponsor compensation and dilution, more information about the target company, and addressing the
use of projections, to align SPACs more closely with traditional IPOs.
\end{dated}

\section{Structural relative value, in general}

The four trades share a shape. A structure ties a price to a value that can be observed or replicated (the other class, the stakes, the net asset value,
the trust). The gap exists because no one can force it closed cheaply: votes, taxes, the inability to redeem, the time to a vote. The trade earns the gap's
narrowing and pays carry while it waits. Its risks are the gap widening (sentiment, as Lee, Shleifer and Thaler found for funds), the hedge failing (a
portfolio that cannot be shorted exactly, as Pontiff found), and time. The practical question is always the same: what closes the gap, and when? With a
catalyst the trade is an event; without one it is a slow, uncertain carry trade.

\section{Strategy files}

\begin{strategyfile}{Share-class spread}\label{strat:s1:share-classes-holding-companies-closed-end-funds-spacs:class}
\sfield{Who pays you, and why} Traders who need one class's liquidity or index membership and push its price away from the other's.
\sfield{Instruments and venues} The two listed classes of one company.
\sfield{Signal} The log spread against its own history, net of the value of the voting difference.
\sfield{Sizing and execution} Long the cheap class, short the rich one; small, frequent trades.
\sfield{Costs} Borrow and spread on the less liquid class.
\sfield{How it dies} Takeover contests that revalue the votes; index changes; conversion or collapse of the structure.
\sfield{Horizon, capacity, infrastructure} Days to weeks; capacity set by the less liquid class.
\sfield{Backtest honestly} Both classes' prices at the same moment; corporate events in the sample.
\sfield{Sources} Chapter~4's twin simulation; no public performance record verified.
\end{strategyfile}

\begin{strategyfile}{Holding-company discount trade}\label{strat:s1:share-classes-holding-companies-closed-end-funds-spacs:holding}
\sfield{Who pays you, and why} Investors who avoid the holding company's structure, taxes or governance, and sell it below its parts.
\sfield{Instruments and venues} The holding company, long; its listed stakes, short.
\sfield{Signal} The discount against its own history and against its causes (tax, costs, governance changes).
\sfield{Sizing and execution} Hedged to the listed stakes; the unlisted part left as risk or estimated.
\sfield{Costs} Borrow on the stakes; carry for years.
\sfield{How it dies} Discounts that persist or widen; no catalyst.
\sfield{Horizon, capacity, infrastructure} Years; the holding's reported portfolio.
\sfield{Backtest honestly} Stakes as reported at each date; the unlisted part marked conservatively.
\sfield{Sources} This chapter's simulation (0.94\% a year without a catalyst); no public performance record verified.
\end{strategyfile}

\begin{strategyfile}{Closed-end fund discount mean reversion}\label{strat:s1:share-classes-holding-companies-closed-end-funds-spacs:cef}
\sfield{Who pays you, and why} Individual investors whose sentiment moves discounts, and who sell funds when discounts are wide.
\sfield{Instruments and venues} Closed-end funds, long; their portfolios or a proxy, short.
\sfield{Signal} The discount against its long-run average and against other funds' discounts.
\sfield{Sizing and execution} Many funds, each hedged; exit near the average discount.
\sfield{Costs} The fund's expenses while held; hedging costs where the portfolio is hard to replicate.
\sfield{How it dies} Sentiment that moves all discounts wider together.
\sfield{Horizon, capacity, infrastructure} Months to years; daily net asset values.
\sfield{Backtest honestly} Net asset values as published (with their lag); expenses and hedge costs.
\sfield{Sources} Lee, Shleifer and Thaler (1991); Pontiff (1996).
\end{strategyfile}

\begin{strategyfile}{Activist-catalyst discount trade}\label{strat:s1:share-classes-holding-companies-closed-end-funds-spacs:activist}
\sfield{Who pays you, and why} Fund managers and boards who keep a discount open until forced to close it.
\sfield{Instruments and venues} Closed-end funds and holding companies with wide discounts and contestable governance.
\sfield{Signal} The discount, and the feasibility of forcing a tender, a liquidation or a conversion.
\sfield{Sizing and execution} Positions large enough to influence the outcome; proxy campaigns.
\sfield{Costs} Legal and campaign costs; long holding periods if the campaign fails.
\sfield{How it dies} Defensive measures; other shareholders who do not support the campaign.
\sfield{Horizon, capacity, infrastructure} One to three years; legal and governance capability.
\sfield{Backtest honestly} Campaigns that failed as well as those that succeeded.
\sfield{Sources} This chapter's simulation (5.0\% a year with 0.3 catalysts a year, against 0.87\% without); no public performance record verified.
\end{strategyfile}

\begin{strategyfile}{SPAC trust-value floor}\label{strat:s1:share-classes-holding-companies-closed-end-funds-spacs:spac}
\sfield{Who pays you, and why} Investors who sell SPAC shares below their \omterm{def:s1:share-classes-holding-companies-closed-end-funds-spacs:spac}{trust value}, and the option value of the merger.
\sfield{Instruments and venues} SPAC shares (and sometimes warrants) before their merger vote.
\sfield{Signal} Price below \omterm{def:s1:share-classes-holding-companies-closed-end-funds-spacs:spac}{trust value} plus accrued interest; time to the vote or deadline.
\sfield{Sizing and execution} Buy below trust, redeem at the vote unless the merged share is worth more; many SPACs.
\sfield{Costs} Low; capital tied up until the vote.
\sfield{How it dies} Trusts earning less than expected; extensions that lengthen the wait; rules and markets that change the product.
\sfield{Horizon, capacity, infrastructure} Months to two years; trust and redemption terms from filings.
\sfield{Backtest honestly} Redemption at \omterm{def:s1:share-classes-holding-companies-closed-end-funds-spacs:spac}{trust value}, not the merged price; the whole SPAC sample.
\sfield{Sources} Klausner, Ohlrogge and Ruan (2022) on net cash and post-merger losses; this chapter's simulation (floor 6.2\%).
\end{strategyfile}

\section{Tutorial: eighty-five cents on the dollar}\label{tut:s1:share-classes-holding-companies-closed-end-funds-spacs:tut}

\begin{tutorial}
\textbf{Goal.} Simulate a discount that reverts slowly and can be closed by a catalyst, trade it, and price a SPAC as a floor and a call. \textbf{End state:}
the table and \cref{fig:s1:share-classes-holding-companies-closed-end-funds-spacs:exit}.
\begin{tutsteps}
\item \textbf{The discount and the trade}: an AR(1) in logs with Poisson catalysts; exit at a level, a catalyst or the horizon.
\omcode{code/firm/structrv/firm_structrv.py}{27}{52}{The discount's paths and the hedged trade.}\label{lst:s1:share-classes-holding-companies-closed-end-funds-spacs:discount}
\item \textbf{The SPAC}: the trust floor and the merger option.
\omcode{code/firm/structrv/firm_structrv.py}{55}{60}{A SPAC share's payoff at the vote.}\label{lst:s1:share-classes-holding-companies-closed-end-funds-spacs:spac}
\item \textbf{Run} \texttt{fund(0.0)}, \texttt{fund(0.1)}, \texttt{fund(0.3)}, \texttt{holding()}, \texttt{spac()} and \texttt{fig\_structrv.py}.
\end{tutsteps}
\textbf{What to change next.} Make the discounts of several funds move together (Lee, Shleifer and Thaler's common sentiment) and size a book of them; add
warrants to the SPAC; let the catalyst's arrival depend on the discount's width.
\end{tutorial}

\section{Build: structural relative value}\label{bld:s1:share-classes-holding-companies-closed-end-funds-spacs:structrv}

\begin{build}
\bfield{Purpose} Discount dynamics with catalysts, the hedged discount trade, and the SPAC floor-plus-option payoff.
\bfield{Interface} \texttt{simulate\_discount(d0, mean, half\_life, vol, days, catalyst\_rate, catalyst\_level, paths, rng)}, \texttt{discount\_trade(paths,
catalyst\_day, close\_level, carry, horizon)}, \texttt{spac\_payoff(trust\_at\_vote, merged\_value)}, \texttt{spac\_return(price, trust, rate, days,
merged\_value)}.
\bfield{Rules} Discounts in logs; carry charged for the time held; exits on the first of level, catalyst or horizon.
\bfield{Acceptance tests} \texttt{code/firm/structrv/tests/}: deterministic reversion with zero volatility; catalysts that close the discount; the trade's exit
and return by hand; the SPAC payoff by hand.
\bfield{Stretch} Correlated discounts across funds; activist campaigns with costs; the SPAC's warrants and extension votes.
\end{build}

\begin{omsources}
\item C.~M.~C.~Lee, A.~Shleifer and R.~H.~Thaler, ``Investor sentiment and the closed-end fund puzzle'', \emph{Journal of Finance} 46(1), 1991.
\item J.~Pontiff, ``Costly arbitrage: evidence from closed-end funds'', \emph{Quarterly Journal of Economics} 111(4), 1996.
\item M.~Klausner, M.~Ohlrogge and E.~Ruan, ``A sober look at SPACs'', \emph{Yale Journal on Regulation} 39, 2022.
\item US SEC, press release 2024-8, SPAC rules, January 2024.
\end{omsources}

\section{Exercises}

\begin{exercise}[$\star$]\label{exo:s1:share-classes-holding-companies-closed-end-funds-spacs:1}
Express discounts of 85, 90 and 95 cents on the dollar in logs.
\end{exercise}

\begin{exercise}[$\star$]\label{exo:s1:share-classes-holding-companies-closed-end-funds-spacs:2}
A trade bought at 85 cents closes at 95 cents after a year with carry of 1.3\%. What is its log return?
\end{exercise}

\begin{exercise}[$\star$]\label{exo:s1:share-classes-holding-companies-closed-end-funds-spacs:3}
A SPAC trades at \$9.80 with a trust of \$10.00 earning 4\% a year and a vote a year away. What is the \omterm{def:s1:share-classes-holding-companies-closed-end-funds-spacs:spac}{trust value} at the vote and the floor return?
\end{exercise}

\begin{exercise}[$\star\star$]\label{exo:s1:share-classes-holding-companies-closed-end-funds-spacs:4}
With a half-life of two years, what fraction of the gap between today's discount and its average remains after one year? Why is the no-catalyst trade so
slow?
\end{exercise}

\begin{exercise}[$\star\star$]\label{exo:s1:share-classes-holding-companies-closed-end-funds-spacs:5}
Why is a \omterm{def:s1:share-classes-holding-companies-closed-end-funds-spacs:cef}{closed-end fund discount} harder to arbitrage than a share-class spread?
\end{exercise}

\begin{exercise}[$\star\star$]\label{exo:s1:share-classes-holding-companies-closed-end-funds-spacs:6}
Explain why a SPAC bought below its \omterm{def:s1:share-classes-holding-companies-closed-end-funds-spacs:spac}{trust value} is a floor plus a call, and what can make the floor fail.
\end{exercise}

\begin{exercise}[$\star\star\star$]\label{exo:s1:share-classes-holding-companies-closed-end-funds-spacs:7}
\emph{Coding.} Run \texttt{fund(0.1)} and compare the annual return and holding time with \texttt{fund(0.0)} and \texttt{fund(0.3)}. How much would you pay a
year for a catalyst arriving at 0.3 a year?
\end{exercise}

\begin{exercise}[$\star\star\star$]\label{exo:s1:share-classes-holding-companies-closed-end-funds-spacs:8}
\emph{Find the flaw.} ``The fund trades at a 15\% discount to its assets; buying it and shorting its portfolio locks in 15\%.''
\end{exercise}

\section{Problem: Eighty-Five Cents on the Dollar}

\begin{problem}[{Weekend problem --- what closes the gap}]\label{pb:s1:share-classes-holding-companies-closed-end-funds-spacs:1}
The chapter's simulated discounts and SPAC, and the public record.

\medskip\noindent\textbf{Part I --- Structures.}
\begin{enumerate}
\item Define a \omterm{def:s1:share-classes-holding-companies-closed-end-funds-spacs:dual}{dual-class share}, a \omterm{def:s1:share-classes-holding-companies-closed-end-funds-spacs:holding}{holding-company discount} and a \omterm{def:s1:share-classes-holding-companies-closed-end-funds-spacs:cef}{closed-end fund discount}.
\item Why do these gaps exist, and why do they persist?
\item Relate share classes to chapter~4's twins.
\item What did Lee, Shleifer and Thaler, and Pontiff, find?
\end{enumerate}

\medskip\noindent\textbf{Part II --- The discount trade.}
\begin{enumerate}[resume]
\item Describe the model of the fund's discount and the trade.
\item Give the results without a catalyst.
\item Give the results with catalysts at 0.1 and 0.3 a year.
\item What does the holding company's result show?
\end{enumerate}

\medskip\noindent\textbf{Part III --- SPACs.}
\begin{enumerate}[resume]
\item Define a SPAC and its \omterm{def:s1:share-classes-holding-companies-closed-end-funds-spacs:spac}{trust value}.
\item Compute the floor for the chapter's SPAC.
\item What did Klausner, Ohlrogge and Ruan find?
\item What does the simulation say about redeeming and holding?
\end{enumerate}

\medskip\noindent\textbf{Part IV --- The verdict.}
\begin{enumerate}[resume]
\item State the \emph{named result}: the expected return and the time to close of a discount trade with and without a catalyst.
\item What do the four structures have in common?
\item Why are discount trades so often activist trades?
\item What did the SEC's 2024 rules change for SPACs?
\item How would you hedge a fund whose portfolio you cannot short?
\item Which strategy file would you run with patient capital?
\item What is the main risk of a book of discount trades?
\item In one sentence: what does structural relative value buy?
\end{enumerate}
\end{problem}

\section{Interview questions}

\begin{interviewq}[$\star$ \iqroles{researcher, trader}]\label{iq:s1:share-classes-holding-companies-closed-end-funds-spacs:1}
Why can a closed-end fund trade below its net asset value for years?
\end{interviewq}

\begin{interviewq}[$\star\star$ \iqroles{trader}]\label{iq:s1:share-classes-holding-companies-closed-end-funds-spacs:2}
How would you trade two share classes of the same company?
\end{interviewq}

\begin{interviewq}[$\star\star$ \iqroles{researcher}]\label{iq:s1:share-classes-holding-companies-closed-end-funds-spacs:3}
How would you value a SPAC share before its merger vote?
\end{interviewq}

\begin{interviewq}[$\star\star$ \iqroles{risk}]\label{iq:s1:share-classes-holding-companies-closed-end-funds-spacs:4}
What are the risks of a portfolio of hedged closed-end fund positions?
\end{interviewq}

\begin{interviewq}[$\star\star$ \iqroles{trader, researcher}]\label{iq:s1:share-classes-holding-companies-closed-end-funds-spacs:5}
A holding company trades at 70\% of its parts. What would make you buy it?
\end{interviewq}

\begin{interviewq}[$\star\star\star$ \iqroles{researcher}]\label{iq:s1:share-classes-holding-companies-closed-end-funds-spacs:6}
A discount follows an Ornstein--Uhlenbeck process toward $\bar d$ with speed $\kappa$, and a catalyst closes it to $d_c$ at Poisson rate $\lambda$. Write the
expected discount at time $t$ and the expected gain of a trade held to $t$.
\end{interviewq}
